% GENERATED FILE. Edit canonical lessons, then run npm run generate:books.
% canonical-source-hash: fnv1a64:8ab4a7638ec0f50f
% canonical-lessons: MR-C548-adthala, MR-C548-todga, MR-C548-pelne, MR-C548-nistarne, MR-C548-dalmalne

\chapter{Obstacles}
\label{ch:mr-obstacles}
\hlchaptermodality{548}

\begin{chapteropening}
\textbf{By the end of this chapter:} \emph{I can talk about obstacles and how people get through them.}

Complete the last lesson of chapter 548: I can talk about obstacles and how people get through them.
\end{chapteropening}

\section[\texorpdfstring{\mr{अडथळा} (aḍthaḷā)}{अडथळा (aḍthaḷā)}]{\mr{अडथळा} (aḍthaḷā) — an obstacle}

\label{lesson:MR-C548-adthala}



\begin{quote}
Before the new one: say the Marathi for a security guard, then the Marathi for to bring under control.
\end{quote}

\subsection*{You'll want to know: \mr{अडथळा}}
\textbf{\mr{अडथळा}} — \emph{aḍthaḷā} — \textquotedblleft{}an obstacle\textquotedblright{}.

\begin{quote}
\mr{रस्त्यात} \mr{पडलेलं} \mr{झाड} \mr{हा} \mr{मोठा} \mr{अडथळा} \mr{होता}.
\end{quote}

\begin{quote}
\emph{rastyāt paḍlelaṁ jhāḍ hā moṭhā aḍthaḷā hotā.}
\end{quote}

\begin{quote}
The tree lying on the road was a big obstacle.
\end{quote}

Masculine.

\begin{cousinweb}
\textbf{\mr{अडथळा}} (\emph{aḍthaḷā}) begins with \mr{अड}, \textquotedblleft{}obstruction\textquotedblright{}: CDIAL lists Marathi \emph{aḍ}, \textquotedblleft{}obstruction\textquotedblright{}, and \emph{aḍṇẽ}, \textquotedblleft{}to be stopped\textquotedblright{}, under \emph{*aḍ}, \textquotedblleft{}obstruct, stop\textquotedblright{} (CDIAL 187). The second part is not settled.
\end{cousinweb}

\begin{grammarlens}[title={What you've built}]
One more word for saying what happens if.
\end{grammarlens}

\subsection*{Your turn}
\begin{itemize}
\raggedright
  \item \emph{Say it:} \emph{aḍthaḷā}
  \item \emph{Say it:} \emph{aḍthaḷā}, once more
  \item \emph{Say it:} \emph{rastyāt paḍlelaṁ jhāḍ hā moṭhā aḍthaḷā hotā.} — The tree lying on the road was a big obstacle.
  \item \emph{Recall:} say the Marathi for a security guard, then the Marathi for to bring under control, then say \emph{aḍthaḷā} again
\end{itemize}

\begin{tcolorbox}[breakable,colback=teal!4,colframe=teal!35!black,title={Before you move on}]
What does \mr{अडथळा} mean? (\textquotedblleft{}An obstacle\textquotedblright{}.) Say it once more.
\end{tcolorbox}
\section[\texorpdfstring{\mr{तोडगा} (toḍgā)}{तोडगा (toḍgā)}]{\mr{तोडगा} (toḍgā) — a way out; a solution}

\label{lesson:MR-C548-todga}



\begin{quote}
Before the new one: say the Marathi for to bring under control, then the Marathi for an obstacle.
\end{quote}

\subsection*{You'll want to know: \mr{तोडगा}}
\textbf{\mr{तोडगा}} — \emph{toḍgā} — \textquotedblleft{}a way out; a solution\textquotedblright{}.

\begin{quote}
\mr{दोन्ही} \mr{भावांनी} \mr{भांडणावर} \mr{एक} \mr{चांगला} \mr{तोडगा} \mr{काढला}.
\end{quote}

\begin{quote}
\emph{donhī bhāvānnī bhāṇḍṇāvar ek cāṅglā toḍgā kāḍhlā.}
\end{quote}

\begin{quote}
The two brothers found a good way out of the quarrel.
\end{quote}

Masculine.

\mr{तोडगा} \mr{काढणे}, to find a way out.

\begin{cousinweb}
\textbf{\mr{तोडगा}} (\emph{toḍgā}) is usually traced to \mr{तोडणे}, \textquotedblleft{}to break\textquotedblright{}, which CDIAL lists under Sanskrit \emph{trōṭayati}, \textquotedblleft{}breaks\textquotedblright{} (CDIAL 6079): a way of breaking a deadlock. The ending -\mr{गा} is not explained there.
\end{cousinweb}

\begin{grammarlens}[title={What you've built}]
One more word for saying what happens if.
\end{grammarlens}

\subsection*{Your turn}
\begin{itemize}
\raggedright
  \item \emph{Say it:} \emph{toḍgā}
  \item \emph{Say it:} \emph{toḍgā}, once more
  \item \emph{Say it:} \emph{donhī bhāvānnī bhāṇḍṇāvar ek cāṅglā toḍgā kāḍhlā.} — The two brothers found a good way out of the quarrel.
  \item \emph{Recall:} say the Marathi for to bring under control, then the Marathi for an obstacle, then say \emph{toḍgā} again
\end{itemize}

\begin{tcolorbox}[breakable,colback=teal!4,colframe=teal!35!black,title={Before you move on}]
What does \mr{तोडगा} mean? (\textquotedblleft{}A way out; a solution\textquotedblright{}.) Say it once more.
\end{tcolorbox}
\section[\texorpdfstring{\mr{पेलणे} (pelṇe)}{पेलणे (pelṇe)}]{\mr{पेलणे} (pelṇe) — to bear a load; to shoulder}

\label{lesson:MR-C548-pelne}



\begin{quote}
Before the new one: say the Marathi for an obstacle, then the Marathi for a way out.
\end{quote}

\subsection*{You'll want to know: \mr{पेलणे}}
\textbf{\mr{पेलणे}} — \emph{pelṇe} — \textquotedblleft{}to bear a load; to shoulder\textquotedblright{}.

\begin{quote}
\mr{घराची} \mr{सगळी} \mr{जबाबदारी} \mr{तिने} \mr{एकटीने} \mr{पेलली}.
\end{quote}

\begin{quote}
\emph{gharācī sagḷī jabābdārī tine ekṭīne pellī.}
\end{quote}

\begin{quote}
She shouldered all the household's responsibility alone.
\end{quote}

Past, with -\mr{ने} (he …): \mr{त्याने} \mr{जबाबदारी} \mr{पेलली}, \emph{tyāne jabābdārī pellī}.

\begin{cousinweb}
\textbf{\mr{पेलणे}} (\emph{pelṇe}) may be inherited: CDIAL lists a Marathi \emph{pelṇẽ}, \textquotedblleft{}to push\textquotedblright{}, under Sanskrit \emph{prērayati}, \textquotedblleft{}sets in motion\textquotedblright{} (CDIAL 9002), but not the sense \textquotedblleft{}to bear a load\textquotedblright{}, so the link is not settled.
\end{cousinweb}

\begin{grammarlens}[title={What you've built}]
One more word for saying what happens if.
\end{grammarlens}

\subsection*{Your turn}
\begin{itemize}
\raggedright
  \item \emph{Say it:} \emph{pelṇe}
  \item \emph{Say it:} \emph{pelṇe}, once more
  \item \emph{Say it:} \emph{gharācī sagḷī jabābdārī tine ekṭīne pellī.} — She shouldered all the household's responsibility alone.
  \item \emph{Recall:} say the Marathi for an obstacle, then the Marathi for a way out, then say \emph{pelṇe} again
\end{itemize}

\begin{tcolorbox}[breakable,colback=teal!4,colframe=teal!35!black,title={Before you move on}]
What does \mr{पेलणे} mean? (\textquotedblleft{}To bear a load; to shoulder\textquotedblright{}.) Say it once more.
\end{tcolorbox}
\section[\texorpdfstring{\mr{निस्तरणे} (nistarṇe)}{निस्तरणे (nistarṇe)}]{\mr{निस्तरणे} (nistarṇe) — to sort out; to set right}

\label{lesson:MR-C548-nistarne}



\begin{quote}
Before the new one: say the Marathi for a way out, then the Marathi for to bear a load.
\end{quote}

\subsection*{You'll want to know: \mr{निस्तरणे}}
\textbf{\mr{निस्तरणे}} — \emph{nistarṇe} — \textquotedblleft{}to sort out; to set right\textquotedblright{}.

\begin{quote}
\mr{भावाने} \mr{केलेली} \mr{चूक} \mr{मला} \mr{निस्तरावी} \mr{लागली}.
\end{quote}

\begin{quote}
\emph{bhāvāne kelelī cūk malā nistarāvī lāglī.}
\end{quote}

\begin{quote}
I had to sort out the mistake my brother made.
\end{quote}

Past, with -\mr{ने} (he …): \mr{त्याने} \mr{चूक} \mr{निस्तरली}, \emph{tyāne cūk nistarlī}.

\begin{cousinweb}
\textbf{\mr{निस्तरणे}} (\emph{nistarṇe}) is usually traced to Sanskrit \emph{nistaraṇa}, \textquotedblleft{}getting across, rescue\textquotedblright{} (Monier-Williams).
\end{cousinweb}

\begin{grammarlens}[title={What you've built}]
One more word for saying what happens if.
\end{grammarlens}

\subsection*{Your turn}
\begin{itemize}
\raggedright
  \item \emph{Say it:} \emph{nistarṇe}
  \item \emph{Say it:} \emph{nistarṇe}, once more
  \item \emph{Say it:} \emph{bhāvāne kelelī cūk malā nistarāvī lāglī.} — I had to sort out the mistake my brother made.
  \item \emph{Recall:} say the Marathi for a way out, then the Marathi for to bear a load, then say \emph{nistarṇe} again
\end{itemize}

\begin{tcolorbox}[breakable,colback=teal!4,colframe=teal!35!black,title={Before you move on}]
What does \mr{निस्तरणे} mean? (\textquotedblleft{}To sort out; to set right\textquotedblright{}.) Say it once more.
\end{tcolorbox}
\section[\texorpdfstring{\mr{डळमळणे} (ḍaḷmaḷṇe)}{डळमळणे (ḍaḷmaḷṇe)}]{\mr{डळमळणे} (ḍaḷmaḷṇe) — to waver; to wobble}

\label{lesson:MR-C548-dalmalne}



\begin{quote}
Before the new one: say the Marathi for to bear a load, then the Marathi for to sort out.
\end{quote}

\subsection*{You'll want to know: \mr{डळमळणे}}
\textbf{\mr{डळमळणे}} — \emph{ḍaḷmaḷṇe} — \textquotedblleft{}to waver; to wobble\textquotedblright{}.

\begin{quote}
\mr{अडचणी} \mr{आल्या} \mr{तरी} \mr{त्याचा} \mr{निर्धार} \mr{डळमळला} \mr{नाही}.
\end{quote}

\begin{quote}
\emph{aḍcaṇī ālyā tarī tyācā nirdhār ḍaḷmaḷlā nāhī.}
\end{quote}

\begin{quote}
Even when difficulties came, his determination did not waver.
\end{quote}

Past (he …): \mr{तो} \mr{डळमळला}, \emph{to ḍaḷmaḷlā}.

\begin{cousinweb}
\textbf{\mr{डळमळणे}} (\emph{ḍaḷmaḷṇe}) is an expressive Marathi verb for wobbling; its origin is not settled.
\end{cousinweb}

\begin{grammarlens}[title={What you've built}]
One more word for saying what happens if.
\end{grammarlens}

\subsection*{Your turn}
\begin{itemize}
\raggedright
  \item \emph{Say it:} \emph{ḍaḷmaḷṇe}
  \item \emph{Say it:} \emph{ḍaḷmaḷṇe}, once more
  \item \emph{Say it:} \emph{aḍcaṇī ālyā tarī tyācā nirdhār ḍaḷmaḷlā nāhī.} — Even when difficulties came, his determination did not waver.
  \item \emph{Recall:} say the Marathi for to bear a load, then the Marathi for to sort out, then say \emph{ḍaḷmaḷṇe} again
\end{itemize}

\begin{tcolorbox}[breakable,colback=teal!4,colframe=teal!35!black,title={Before you move on}]
What does \mr{डळमळणे} mean? (\textquotedblleft{}To waver; to wobble\textquotedblright{}.) Say it once more.
\end{tcolorbox}
